% SKELETON: This file contains placeholder result tables (---).
% See paper.tex for the complete version with actual experimental results.
\documentclass{article}

% NeurIPS 2024 style
\usepackage[preprint]{neurips_2024}

\usepackage[utf8]{inputenc}
\usepackage[T1]{fontenc}
\usepackage{hyperref}
\usepackage{url}
\usepackage{booktabs}
\usepackage{amsfonts}
\usepackage{amsmath}
\usepackage{amssymb}
\usepackage{nicefrac}
\usepackage{microtype}
\usepackage{xcolor}
\usepackage{graphicx}
\usepackage{multirow}
\usepackage{array}
\usepackage{tabularx}
\usepackage{algorithm}
\usepackage{algpseudocode}
\usepackage{subcaption}

\title{Obstacle-Aware Locomotion for Humanoid Robots:\\
       A Comparative Study of PPO and TD3 with LiDAR Perception}

\author{
  Parth \\
  Project Saraswati \\
  \texttt{unitree\_rl\_lab}
}

\begin{document}

\maketitle

% ─────────────────────────────────────────────────────────────
\begin{abstract}
% ─────────────────────────────────────────────────────────────
We study the problem of obstacle-aware locomotion for the Unitree G1 29-DoF
humanoid robot using deep reinforcement learning in simulation.
We formulate two distinct Markov Decision Processes: an \emph{empty world} for
baseline locomotion and an \emph{obstacle world} that introduces randomly
placed cuboid obstacles with collision termination and proximity penalties.
For each world we train two families of agents-one that relies solely on
proprioceptive observations and one that additionally receives a front-facing
LiDAR point cloud-using both Proximal Policy Optimization (PPO) and Twin
Delayed Deep Deterministic Policy Gradient (TD3) as learning algorithms.
Through systematic ablations we analyse (i)~the effect of the learning
algorithm in isolation (PPO vs.\ TD3), (ii)~the contribution of LiDAR
perception, and (iii)~the sim-to-sim transfer gap between agents trained in
empty and obstacle environments.
All experiments are conducted in NVIDIA Isaac Lab using 4{,}096 parallel
environments on a single GPU, enabling rapid iteration while maintaining
physical fidelity.
\end{abstract}

% ─────────────────────────────────────────────────────────────
\section{Introduction}
% ─────────────────────────────────────────────────────────────

Bipedal locomotion in cluttered environments remains a central challenge in
robotics.
While model-based controllers excel in structured settings, they degrade
gracefully under sensor noise and unmodelled contacts.
Reinforcement learning (RL)-particularly when coupled with massively parallel
physics simulation-has recently produced locomotion controllers that match or
exceed model-based alternatives on rough terrain
\cite{lee2020learning,kumar2021rma,margolis2022rapid}.
However, most prior work focuses on terrain adaptation and ignores dynamic or
static obstacles that require active avoidance.

The Unitree G1, with its 29 degrees of freedom and human-proportioned
morphology, presents a rich testbed for obstacle-aware locomotion research.
We make the following contributions:

\begin{itemize}
  \item A sim-ready MDP formulation for a 29-DoF humanoid in both free and
        obstacle-cluttered environments, implemented in Isaac Lab.
  \item A systematic ablation comparing on-policy (PPO) and off-policy (TD3)
        algorithms across two observation modalities (proprioception only vs.\
        proprioception + LiDAR).
  \item An empirical analysis of curriculum-free transfer from empty-world to
        obstacle-world policies.
  \item All code, configurations, and training scripts released with the paper.
\end{itemize}

% ─────────────────────────────────────────────────────────────
\section{Related Work}
% ─────────────────────────────────────────────────────────────

\paragraph{Legged locomotion via RL.}
\citet{lee2020learning} demonstrated that a teacher–student framework can
transfer rough-terrain locomotion policies from simulation to real quadrupeds.
\citet{kumar2021rma} introduced Rapid Motor Adaptation (RMA), learning an
online adaptation module conditioned on extrinsics.
For bipeds, \citet{radosavovic2024humanoid} and concurrent work show
sim-to-real transfer of whole-body locomotion on the Unitree H1 platform.

\paragraph{Obstacle avoidance for legged robots.}
Navigation-aware locomotion has been studied by coupling high-level planners
with low-level locomotion policies
\cite{hoeller2024anymal,zhuang2023robot}.
Our work differs in that we embed obstacle awareness directly into the
locomotion policy via LiDAR observations rather than relying on a separate
planning layer.

\paragraph{On-policy vs.\ off-policy RL for continuous control.}
PPO \cite{schulman2017proximal} dominates locomotion benchmarks due to
sample efficiency in the massively-parallel setting.
TD3 \cite{fujimoto2018addressing} has demonstrated strong results in
lower-dimensional MuJoCo tasks but is less studied at humanoid scale.
We provide the first direct comparison of the two on a 29-DoF humanoid.

% ─────────────────────────────────────────────────────────────
\section{Background}
% ─────────────────────────────────────────────────────────────

\subsection{Markov Decision Process}

We formulate each task as a finite-horizon MDP
$\mathcal{M} = (\mathcal{S}, \mathcal{A}, \mathcal{T}, r, \gamma, \rho_0)$,
where $\mathcal{S}$ is the state space, $\mathcal{A}$ the action space,
$\mathcal{T}$ the (unknown) transition dynamics, $r$ the reward function,
$\gamma \in [0,1)$ the discount factor, and $\rho_0$ the initial state
distribution.
The agent observes a partial observation $o_t \in \mathcal{O}$ and learns a
policy $\pi_\theta : \mathcal{O} \to \mathcal{A}$.

\subsection{Proximal Policy Optimization (PPO)}

PPO \cite{schulman2017proximal} is an on-policy actor-critic algorithm that
constrains policy updates via a clipped surrogate objective:
\begin{equation}
  \mathcal{L}^{\text{CLIP}}(\theta)
  = \mathbb{E}_t\!\left[
      \min\!\left(
        r_t(\theta)\,\hat{A}_t,\;
        \operatorname{clip}\!\left(r_t(\theta),\,1-\epsilon,\,1+\epsilon\right)\hat{A}_t
      \right)
    \right],
\end{equation}
where $r_t(\theta) = \nicefrac{\pi_\theta(a_t|o_t)}{\pi_{\theta_{\text{old}}}(a_t|o_t)}$
is the probability ratio and $\hat{A}_t$ is the Generalised Advantage
Estimate (GAE) \cite{schulman2015high}.
The full objective adds a value-function loss and an entropy bonus:
\begin{equation}
  \mathcal{L}(\theta)
  = \mathcal{L}^{\text{CLIP}}(\theta)
    - c_1 \mathcal{L}^{\text{VF}}(\theta)
    + c_2 \mathcal{H}[\pi_\theta].
\end{equation}
PPO is particularly well-suited to massively-parallel simulation because
rollouts from all environments can be batched into a single update step.

\subsection{Twin Delayed Deep Deterministic Policy Gradient (TD3)}

TD3 \cite{fujimoto2018addressing} is an off-policy actor-critic algorithm
that extends DDPG \cite{lillicrap2015continuous} with three key
improvements: (i)~\emph{clipped double-Q learning} to reduce overestimation
bias, (ii)~\emph{delayed policy updates} (the actor is updated every $d$
critic steps), and (iii)~\emph{target policy smoothing} via noise injected
into the target action:
\begin{align}
  y &= r + \gamma\,\min_{i=1,2} Q_{\theta'_i}\!\left(s', \pi_{\phi'}(s') + \epsilon\right),
      \quad \epsilon \sim \operatorname{clip}(\mathcal{N}(0,\sigma),{-c},{c}), \\
  \mathcal{L}(\theta_i) &= \mathbb{E}\!\left[(y - Q_{\theta_i}(s,a))^2\right].
\end{align}
TD3 learns from a replay buffer $\mathcal{D}$ and is therefore sample-efficient
for long-horizon tasks, though it requires careful tuning of exploration noise.

% ─────────────────────────────────────────────────────────────
\section{MDP Formulation}
% ─────────────────────────────────────────────────────────────

\subsection{Two-World Curriculum}

We define two distinct environments:

\begin{description}
  \item[Empty World ($\mathcal{E}$).] A flat infinite ground plane with no
        obstacles. The robot is rewarded purely for velocity tracking and
        locomotion quality. This world is used to establish a locomotion
        baseline and pre-train policies that are later fine-tuned.
  \item[Obstacle World ($\mathcal{W}$).] The same ground plane augmented with
        up to six kinematic cuboid obstacles, randomly placed at each episode
        reset. Obstacle positions, counts, and the corridor in front of the
        robot are all randomised to prevent overfitting to a fixed layout.
\end{description}

\subsection{Robot Platform}

We use the \textbf{Unitree G1} humanoid robot with \textbf{29 actuated
joints}: 12 for the legs, 7 per arm, and 3 for the waist.
The robot stands approximately 1.3\,m tall with a target standing height of
0.78\,m for the base link.
All joints are controlled via \emph{position targets} with a PD controller
running at 200\,Hz; the RL policy runs at 50\,Hz (decimation factor 4).

\subsection{Observation Space}

We study two observation modalities.

\paragraph{Without LiDAR (proprioceptive only).}
The policy observation $o_t$ at each step contains:

\begin{center}
\begin{tabular}{lcc}
  \toprule
  \textbf{Component} & \textbf{Dim} & \textbf{Scale / Noise} \\
  \midrule
  Base angular velocity & 3 & $\times 0.2$, $\sigma{=}0.2$ \\
  Projected gravity vector & 3 & $\sigma{=}0.05$ \\
  Velocity commands $(v_x, v_y, \omega_z)$ & 3 & --- \\
  Joint positions (rel.\ default) & 29 & $\sigma{=}0.01$ \\
  Joint velocities & 29 & $\times 0.05$, $\sigma{=}1.5$ \\
  Last action & 29 & --- \\
  \midrule
  \textbf{Per-step total} & \textbf{96} & \\
  \textbf{With history ($T{=}5$)} & \textbf{480} & \\
  \bottomrule
\end{tabular}
\end{center}

The critic additionally receives the base linear velocity (3D), giving a
483-dimensional critic observation per step.
Gaussian observation noise is applied during training only to improve
robustness.

\paragraph{With LiDAR.}
A \texttt{MultiMeshRayCaster} sensor is mounted on the torso link at offset
$(0.2, 0.0, 0.15)$\,m and aligned to the robot's base frame.
It fires $8\,\text{channels} \times 31\,\text{rays} = 248$ rays over a
$180^\circ$ horizontal and $30^\circ$ vertical field of view at 6° resolution
with a maximum range of 20\,m.
LiDAR ranges are appended to both policy and critic observations, increasing
the dimension by $248 \times 5 = 1{,}240$ per history window.
No explicit normalization is applied to the LiDAR; empirical normalization
inside the PPO runner handles scaling automatically.

\subsection{Action Space}

The action $a_t \in \mathbb{R}^{29}$ represents \emph{joint position offsets}
from the robot's default standing pose, scaled by 0.25 radians.
A low-level PD controller converts these targets to joint torques within the
simulator.

\subsection{Reward Functions}

\subsubsection{Empty World Rewards ($n-1$ terms)}

The empty-world reward $r_t^\mathcal{E}$ consists of the following 19 terms:

\begin{center}
\begin{tabular}{llr}
  \toprule
  \textbf{Term} & \textbf{Description} & \textbf{Weight} \\
  \midrule
  \multicolumn{3}{l}{\textit{Tracking}} \\
  \quad\texttt{track\_lin\_vel\_xy} & Exp.\ tracking of commanded $v_x, v_y$ & $+1.0$ \\
  \quad\texttt{track\_ang\_vel\_z} & Exp.\ tracking of commanded $\omega_z$ & $+0.5$ \\
  \midrule
  \multicolumn{3}{l}{\textit{Survival}} \\
  \quad\texttt{alive} & Constant bonus per step & $+0.15$ \\
  \midrule
  \multicolumn{3}{l}{\textit{Smoothness \& Efficiency}} \\
  \quad\texttt{base\_linear\_velocity} & Penalise $z$-axis linear velocity & $-2.0$ \\
  \quad\texttt{base\_angular\_velocity} & Penalise $xy$ angular velocity & $-0.05$ \\
  \quad\texttt{joint\_vel} & L2 joint velocity & $-0.001$ \\
  \quad\texttt{joint\_acc} & L2 joint acceleration & $-2.5\times10^{-7}$ \\
  \quad\texttt{action\_rate} & Penalise action change rate & $-0.05$ \\
  \quad\texttt{energy} & Penalise $\tau \cdot \dot{q}$ & $-2\times10^{-5}$ \\
  \midrule
  \multicolumn{3}{l}{\textit{Posture}} \\
  \quad\texttt{flat\_orientation\_l2} & Penalise body tilt & $-5.0$ \\
  \quad\texttt{base\_height} & Penalise deviation from 0.78\,m & $-10.0$ \\
  \quad\texttt{dof\_pos\_limits} & Joint limit violation & $-5.0$ \\
  \quad\texttt{joint\_deviation\_arms} & Arm joints near default & $-0.1$ \\
  \quad\texttt{joint\_deviation\_waists} & Waist joints near default & $-1.0$ \\
  \quad\texttt{joint\_deviation\_legs} & Hip roll/yaw near default & $-1.0$ \\
  \midrule
  \multicolumn{3}{l}{\textit{Gait Quality}} \\
  \quad\texttt{gait} & Phase-synchronised footstep reward & $+0.5$ \\
  \quad\texttt{feet\_clearance} & Swing-foot height clearance & $+1.0$ \\
  \quad\texttt{feet\_slide} & Penalise foot sliding & $-0.2$ \\
  \quad\texttt{undesired\_contacts} & Penalise non-foot contacts & $-1.0$ \\
  \bottomrule
\end{tabular}
\end{center}

The velocity-tracking terms use the exponential kernel
$r = \exp(-\|e\|^2 / \sigma^2)$ with $\sigma^2 = 0.25$.
The gait reward enforces a walking pattern with period $T_g = 0.8$\,s and
phase offset $\Delta\phi = 0.5$ (alternating feet), using a stance threshold
of 0.55.

\subsubsection{Obstacle World Rewards ($n$ terms)}

The obstacle-world reward $r_t^\mathcal{W}$ extends $r_t^\mathcal{E}$ with
two additional terms:

\begin{center}
\begin{tabular}{llr}
  \toprule
  \textbf{Term} & \textbf{Description} & \textbf{Weight} \\
  \midrule
  \texttt{object\_hit} & Torso contact force with obstacles & $-2.0$ \\
  \texttt{lidar\_proximity} & Rays within $d_\text{safe}{=}0.8$\,m of obstacle & $-0.5$ \\
  \bottomrule
\end{tabular}
\end{center}

The \texttt{object\_hit} penalty is computed from the filtered contact force
matrix $\|\mathbf{F}_\text{mat}\|$ (obstacles only, via
\texttt{filter\_prim\_paths\_expr}), ensuring ground contacts do not
contribute.
The \texttt{lidar\_proximity} penalty sums
$\max(0,\, d_\text{safe} - d_i) / d_\text{safe}$ over all 248 rays,
providing a smooth keep-away signal before physical contact occurs.

\subsection{Termination Conditions}

\begin{center}
\begin{tabular}{ll}
  \toprule
  \textbf{Condition} & \textbf{Threshold} \\
  \midrule
  \texttt{time\_out} & Episode length $> 20$\,s \\
  \texttt{base\_height} & Root height $< 0.2$\,m \\
  \texttt{bad\_orientation} & Body tilt $> 0.8$\,rad \\
  \texttt{obstacle\_collision}$^*$ & Contact force with obstacle $> 1$\,N \\
  \bottomrule
  \multicolumn{2}{l}{$^*$Obstacle world only.}
\end{tabular}
\end{center}

% ─────────────────────────────────────────────────────────────
\section{Algorithms and Hyperparameters}
% ─────────────────────────────────────────────────────────────

\subsection{PPO Configuration}

We use the RSL-RL implementation of PPO \cite{rudin2022learning} with an
asymmetric actor-critic architecture.
The actor receives the policy observation group; the critic receives a
privileged observation group that additionally includes the ground-truth base
linear velocity.

\begin{center}
\begin{tabular}{ll}
  \toprule
  \textbf{Hyperparameter} & \textbf{Value} \\
  \midrule
  \multicolumn{2}{l}{\textit{Network architecture (both actor and critic)}} \\
  \quad Hidden layers & [1024, 1024, 512, 256, 128] \\
  \quad Activation & ELU \\
  \quad Output distribution & Gaussian, $\sigma_0 = 1.0$ \\
  \midrule
  \multicolumn{2}{l}{\textit{Rollout \& update}} \\
  \quad Steps per env per update & 24 \\
  \quad Number of envs & 4{,}096 \\
  \quad Mini-batches & 4 \\
  \quad Learning epochs per update & 5 \\
  \quad Total rollout steps (batch) & 98{,}304 \\
  \midrule
  \multicolumn{2}{l}{\textit{Optimisation}} \\
  \quad Learning rate & $1\times10^{-3}$ (adaptive) \\
  \quad Clip parameter $\epsilon$ & 0.2 \\
  \quad Entropy coefficient $c_2$ & 0.01 \\
  \quad Value loss coefficient $c_1$ & 1.0 \\
  \quad Discount factor $\gamma$ & 0.99 \\
  \quad GAE parameter $\lambda$ & 0.95 \\
  \quad Desired KL divergence & 0.01 \\
  \quad Max gradient norm & 1.0 \\
  \quad Empirical normalization & True \\
  \midrule
  \multicolumn{2}{l}{\textit{Training}} \\
  \quad Max iterations & 50{,}000 \\
  \quad Best checkpoint criterion & Max mean episode return \\
  \bottomrule
\end{tabular}
\end{center}

The adaptive learning rate schedule adjusts $\alpha$ up or down by 1.5$\times$
based on whether the measured KL divergence exceeds the desired threshold,
stabilising training under varying reward scales.

\subsection{TD3 Configuration}

We use the Stable-Baselines3 implementation of TD3 \cite{raffin2021stable}
with the same MLP architecture as PPO to ensure a fair comparison of the
algorithmic components.

\begin{center}
\begin{tabular}{ll}
  \toprule
  \textbf{Hyperparameter} & \textbf{Value} \\
  \midrule
  \multicolumn{2}{l}{\textit{Network architecture (policy \& Q-networks)}} \\
  \quad Hidden layers & [1024, 1024, 512, 256, 128] \\
  \quad Activation & ReLU \\
  \quad Twin Q-networks & Yes \\
  \midrule
  \multicolumn{2}{l}{\textit{Replay buffer \& sampling}} \\
  \quad Replay buffer size & 1{,}000{,}000 \\
  \quad Batch size & 256 \\
  \quad Learning starts & 10{,}000 steps \\
  \quad Train frequency & 1 step \\
  \quad Gradient steps per step & 1 \\
  \midrule
  \multicolumn{2}{l}{\textit{Optimisation}} \\
  \quad Learning rate & $3\times10^{-4}$ \\
  \quad Discount factor $\gamma$ & 0.99 \\
  \quad Soft update coefficient $\tau$ & 0.005 \\
  \quad Policy delay $d$ & 2 \\
  \quad Target policy noise $\sigma$ & 0.2 \\
  \quad Noise clip $c$ & 0.5 \\
  \midrule
  \multicolumn{2}{l}{\textit{Training}} \\
  \quad Total timesteps & 1{,}000{,}000 \\
  \quad Checkpoint frequency & 50{,}000 steps \\
  \quad Seed & 42 \\
  \bottomrule
\end{tabular}
\end{center}

The policy delay of $d=2$ means the actor is updated once for every two
critic updates, reducing variance in policy gradient estimates.
Target policy smoothing with $\sigma=0.2$ acts as a regulariser that prevents
the policy from exploiting narrow peaks in the Q-function.

% ─────────────────────────────────────────────────────────────
\section{Experimental Setup}
% ─────────────────────────────────────────────────────────────

\subsection{Simulation Framework}

All experiments are conducted in \textbf{NVIDIA Isaac Lab} \cite{mittal2023orbit},
a modular framework built on top of Isaac Sim (PhysX 5) that provides
GPU-accelerated parallel simulation.
We run 4{,}096 parallel environments on a single NVIDIA GPU, achieving
approximately 50{,}000 environment steps per second for proprioceptive
policies.

\begin{center}
\begin{tabular}{ll}
  \toprule
  \textbf{Simulation Parameter} & \textbf{Value} \\
  \midrule
  Physics engine & PhysX 5 (GPU) \\
  Physics timestep $\Delta t_\text{sim}$ & 5\,ms (200\,Hz) \\
  Control decimation & $4\times$ (50\,Hz) \\
  Number of training environments & 4{,}096 \\
  Environment spacing (empty world) & 2.5\,m \\
  Environment spacing (obstacle world) & 6.0\,m \\
  Ground plane & Infinite (plane), $\mu_s{=}1.0$ \\
  \midrule
  \multicolumn{2}{l}{\textit{Obstacle world specific}} \\
  Obstacle count per env & 6 cuboids (0–5 active per reset) \\
  Obstacle dimensions & $(0.3\text{–}0.55) \times (0.3\text{–}0.55) \times (1.8\text{–}3.4)$\,m \\
  Placement range $x$ & [1.5, 5.0]\,m ahead of robot \\
  Placement range $y$ & $[-2.5, 2.5]$\,m lateral \\
  Minimum obstacle separation & 1.25\,m \\
  \bottomrule
\end{tabular}
\end{center}

\subsection{Obstacle Randomisation}

At every episode reset, the event manager samples a uniformly random number of
active obstacles $N_\text{active}$ from a curriculum-dependent range
(Table~\ref{tab:curricula}) and places them using rejection sampling to
enforce minimum separation and a spawn-clearance zone around the robot.
Inactive obstacles are teleported underground at $z = -10$\,m.
The LiDAR raycasts against all obstacle prims across all environments using
a global path expression, enabling each robot to perceive obstacles from
neighbouring environments as well as its own-a form of implicit data
augmentation.

\begin{table}[h]
  \centering
  \caption{Obstacle environment variants.}
  \label{tab:curricula}
  \begin{tabular}{lcc}
    \toprule
    \textbf{Variant} & $N_\text{min}$ & $N_\text{max}$ \\
    \midrule
    Single obstacle & 1 & 1 \\
    Random obstacle & 0 & 3 \\
    Dense obstacle  & 3 & 5 \\
    \bottomrule
  \end{tabular}
\end{table}

\subsection{Training Protocol}

For PPO we adopt a two-stage training strategy:
\begin{enumerate}
  \item \textbf{Stage 1 – Empty world pre-training.} Train for up to 50{,}000
        iterations in $\mathcal{E}$ until the policy achieves stable
        forward locomotion.
  \item \textbf{Stage 2 – Obstacle world fine-tuning.} Initialise from the
        Stage~1 checkpoint and continue training in $\mathcal{W}$ with the
        obstacle reward terms and collision termination enabled.
\end{enumerate}
For TD3 we train directly in each world from scratch due to the off-policy
nature of the algorithm.
All training runs log the mean episode return across environments and
checkpoint both the latest best model (\texttt{best\_model.pt}) and an
iteration-stamped snapshot (\texttt{best\_model\_\{iter\}.pt}) whenever a
new best return is achieved.

% ─────────────────────────────────────────────────────────────
\section{Results}
% ─────────────────────────────────────────────────────────────

We report results across six ablation conditions.
For each condition we record the mean episode return (MER), mean episode
length (MEL), and the collision rate (CR)-the fraction of episodes that
terminate due to obstacle contact-where applicable.

\subsection{Ablation I: PPO vs.\ TD3 - Empty World, No LiDAR}

\begin{center}
\begin{tabular}{lcccc}
  \toprule
  \textbf{Algorithm} & \textbf{MER} $\uparrow$ & \textbf{MEL (s)} $\uparrow$
    & \textbf{Iterations / Steps} \\
  \midrule
  PPO (ours) & --- & --- & 50{,}000 iters \\
  TD3         & --- & --- & 1M steps \\
  \bottomrule
\end{tabular}
\end{center}

\textit{[Fill in numerical results after training runs complete.]}

The massively-parallel rollout collection of PPO provides a significant
wall-clock advantage in this setting: 4{,}096 environments collecting 24
steps each yield $\sim$100K transitions per update, whereas TD3 collects
transitions serially from a single environment instance and must warm-start
its replay buffer before useful gradient updates begin.
We therefore expect PPO to reach locomotion competence substantially faster
in wall-clock time, despite potentially lower asymptotic sample efficiency.

\subsection{Ablation II: PPO vs.\ TD3 - Obstacle World, No LiDAR}

\begin{center}
\begin{tabular}{lcccc}
  \toprule
  \textbf{Algorithm} & \textbf{MER} $\uparrow$ & \textbf{MEL (s)} $\uparrow$
    & \textbf{CR} $\downarrow$ \\
  \midrule
  PPO (ours) & --- & --- & --- \\
  TD3         & --- & --- & --- \\
  \bottomrule
\end{tabular}
\end{center}

\textit{[Fill in numerical results after training runs complete.]}

Without LiDAR the robot must rely on reactive proprioception to avoid
obstacles after physical contact has already begun (signalled by the
\texttt{object\_hit} reward).
We hypothesise that PPO's advantage in this regime will narrow because the
on-policy rollouts frequently contain collision events early in training,
leading to highly truncated episodes and noisy advantage estimates.

\subsection{Ablation III \& IV: Algorithm Comparison - With LiDAR}

\begin{center}
\begin{tabular}{llccc}
  \toprule
  \textbf{World} & \textbf{Algorithm} & \textbf{MER} $\uparrow$
    & \textbf{MEL (s)} $\uparrow$ & \textbf{CR} $\downarrow$ \\
  \midrule
  Empty    & PPO + LiDAR & --- & --- & N/A \\
  Empty    & TD3 + LiDAR & --- & --- & N/A \\
  Obstacle & PPO + LiDAR & --- & --- & --- \\
  Obstacle & TD3 + LiDAR & --- & --- & --- \\
  \bottomrule
\end{tabular}
\end{center}

\textit{[Fill in numerical results after training runs complete.]}

The LiDAR observation adds 1{,}240 dimensions to the policy input.
For PPO this is handled transparently by the existing MLP; the
\texttt{empirical\_normalization} module accumulates running statistics to
scale the additional inputs.
For TD3 the larger observation space increases the input dimension of both
the policy and Q-networks, which may require additional replay buffer
warm-up.

\subsection{Ablation V: Transfer - Empty-World Policy in Obstacle World}

\begin{center}
\begin{tabular}{lccc}
  \toprule
  \textbf{Policy trained in} & \textbf{Evaluated in} & \textbf{MER} $\uparrow$
    & \textbf{CR} $\downarrow$ \\
  \midrule
  Empty world (no LiDAR) & Obstacle world & --- & --- \\
  Obstacle world (no LiDAR) & Obstacle world & --- & --- \\
  Empty world (LiDAR) & Obstacle world & --- & --- \\
  Obstacle world (LiDAR) & Obstacle world & --- & --- \\
  \bottomrule
\end{tabular}
\end{center}

\textit{[Fill in numerical results after training runs complete.]}

This ablation isolates the \emph{locomotion quality gap}: a policy trained
without obstacles but with a well-shaped gait should navigate simple obstacle
configurations through reactive avoidance, while a policy trained in the
obstacle world should exhibit more deliberate path planning.

% ─────────────────────────────────────────────────────────────
\section{Results Discussion}
% ─────────────────────────────────────────────────────────────

\paragraph{PPO vs.\ TD3.}
The primary expected finding is that PPO converges faster in wall-clock time
due to massively-parallel rollout collection, but TD3 may achieve a higher
asymptotic performance given its off-policy replay buffer that allows each
transition to be reused many times.
In the obstacle world, the collision termination introduces highly variable
episode lengths which can destabilise advantage estimates for PPO; TD3's
replay buffer naturally averages over this variability.

\paragraph{Effect of LiDAR.}
LiDAR is expected to be most beneficial in the obstacle world where
anticipatory avoidance (before contact) is rewarded by the
\texttt{lidar\_proximity} term.
In the empty world, the additional 1{,}240-dimensional input may moderately
hurt convergence speed without providing a meaningful performance improvement.

\paragraph{Empty-world pre-training.}
Two-stage training (pre-train on empty, fine-tune on obstacles) is expected
to outperform training from scratch in the obstacle world because the robot
first acquires stable walking, which is a prerequisite for effective obstacle
avoidance.
Transfer without fine-tuning is expected to produce high collision rates due
to the absence of any obstacle-aware reward signal during training.

\paragraph{Limitations.}
All results are obtained in simulation with kinematic (static) obstacles.
Generalisation to dynamic obstacles and real-world deployment remains future
work.
The TD3 baseline uses a single environment for experience collection; a
multi-environment wrapper for TD3 would provide a fairer comparison.

% ─────────────────────────────────────────────────────────────
\section{Future Work}
% ─────────────────────────────────────────────────────────────

\begin{itemize}
  \item \textbf{Dynamic obstacles.} Extending the obstacle world to include
        moving objects (other robots, pedestrians) and evaluating whether the
        LiDAR policy generalises without retraining.
  \item \textbf{Sim-to-real transfer.} Deploying the best obstacle-aware
        policy on the physical Unitree G1 using domain randomisation over
        physics parameters and sensor noise.
  \item \textbf{Curriculum learning.} Automating the progression from empty
        to dense-obstacle environments via adaptive difficulty, e.g., PAIRED
        \cite{dennis2020emergent} or automatic curriculum generation based on
        episode return.
  \item \textbf{Multi-modal fusion.} Combining LiDAR with onboard camera data
        for richer scene understanding, potentially using a learned encoder
        rather than raw range values.
  \item \textbf{Whole-body obstacle avoidance.} Extending the contact sensor
        to all body links (not just the torso) and adding per-body proximity
        penalties to encourage avoidance with arms, head, and legs.
  \item \textbf{Multi-environment TD3.} Implementing a vectorised replay
        buffer for TD3 to leverage the same 4{,}096-environment
        infrastructure as PPO, enabling a truly fair comparison.
\end{itemize}

% ─────────────────────────────────────────────────────────────
\section*{Acknowledgements}
% ─────────────────────────────────────────────────────────────

We thank the Isaac Lab and RSL-RL open-source communities for the simulation
infrastructure and PPO implementation used in this work.

% ─────────────────────────────────────────────────────────────
\bibliographystyle{plainnat}
\begin{thebibliography}{99}

\bibitem[Dennis et al.(2020)]{dennis2020emergent}
Dennis, M., Jaques, N., Vinitsky, E., Bayen, A., Russell, S., Critch, A.,
and Levine, S.
\newblock Emergent complexity and zero-shot transfer via unsupervised
environment design.
\newblock \emph{Advances in Neural Information Processing Systems (NeurIPS)},
2020.

\bibitem[Fujimoto et al.(2018)]{fujimoto2018addressing}
Fujimoto, S., van Hoof, H., and Meger, D.
\newblock Addressing function approximation error in actor-critic methods.
\newblock \emph{International Conference on Machine Learning (ICML)}, 2018.

\bibitem[Hoeller et al.(2024)]{hoeller2024anymal}
Hoeller, D., Rudin, N., Sako, D., and Hutter, M.
\newblock Anymal parkour: Learning agile navigation for quadrupedal robots.
\newblock \emph{Science Robotics}, 2024.

\bibitem[Kumar et al.(2021)]{kumar2021rma}
Kumar, A., Fu, Z., Pathak, D., and Malik, J.
\newblock RMA: Rapid motor adaptation for legged robots.
\newblock \emph{Robotics: Science and Systems (RSS)}, 2021.

\bibitem[Lee et al.(2020)]{lee2020learning}
Lee, J., Hwangbo, J., Wellhausen, L., Koltun, V., and Hutter, M.
\newblock Learning quadrupedal locomotion over challenging terrain.
\newblock \emph{Science Robotics}, 5(47), 2020.

\bibitem[Lillicrap et al.(2015)]{lillicrap2015continuous}
Lillicrap, T.~P., Hunt, J.~J., Pritzel, A., Heess, N., Erez, T., Tassa, Y.,
Silver, D., and Wierstra, D.
\newblock Continuous control with deep reinforcement learning.
\newblock \emph{International Conference on Learning Representations (ICLR)},
2016.

\bibitem[Margolis et al.(2022)]{margolis2022rapid}
Margolis, G.~B., Yang, G., Paigwar, K., Chen, T., and Agrawal, P.
\newblock Rapid locomotion via reinforcement learning.
\newblock \emph{Robotics: Science and Systems (RSS)}, 2022.

\bibitem[Mittal et al.(2023)]{mittal2023orbit}
Mittal, M., Yu, C., Yu, Q., Liu, J., Rudin, N., Hoeller, D., Yuan, J.~L.,
Singh, R., Guo, Y., Mazhar, H., Mandlekar, A., Babich, B., State, G.,
Hutter, M., and Garg, A.
\newblock Orbit: A unified simulation framework for interactive robot learning
environments.
\newblock \emph{IEEE Robotics and Automation Letters}, 8(6), 2023.

\bibitem[Radosavovic et al.(2024)]{radosavovic2024humanoid}
Radosavovic, I., Shi, B., Fu, L., Gao, K., Goldberg, K., and Malik, J.
\newblock Humanoid locomotion as next token prediction.
\newblock \emph{arXiv preprint arXiv:2402.19469}, 2024.

\bibitem[Raffin et al.(2021)]{raffin2021stable}
Raffin, A., Hill, A., Gleave, A., Kanervisto, A., Ernestus, M., and
Dormann, N.
\newblock Stable-Baselines3: Reliable reinforcement learning implementations.
\newblock \emph{Journal of Machine Learning Research}, 22(268):1–8, 2021.

\bibitem[Rudin et al.(2022)]{rudin2022learning}
Rudin, N., Hoeller, D., Reist, P., and Hutter, M.
\newblock Learning to walk in minutes using massively parallel deep
reinforcement learning.
\newblock \emph{Conference on Robot Learning (CoRL)}, 2022.

\bibitem[Schulman et al.(2015)]{schulman2015high}
Schulman, J., Moritz, P., Levine, S., Jordan, M., and Abbeel, P.
\newblock High-dimensional continuous control using generalized advantage
estimation.
\newblock \emph{International Conference on Learning Representations (ICLR)},
2016.

\bibitem[Schulman et al.(2017)]{schulman2017proximal}
Schulman, J., Wolski, F., Dhariwal, P., Radford, A., and Klimov, O.
\newblock Proximal policy optimization algorithms.
\newblock \emph{arXiv preprint arXiv:1707.06347}, 2017.

\bibitem[Zhuang et al.(2023)]{zhuang2023robot}
Zhuang, Z., Fu, Z., Wang, J., Atkeson, C., Schwertfeger, S., Finn, C.,
and Zhao, H.
\newblock Robot parkour learning.
\newblock \emph{Conference on Robot Learning (CoRL)}, 2023.

\end{thebibliography}

\appendix

% ─────────────────────────────────────────────────────────────
\section{Reward Function Definitions}
\label{app:rewards}
% ─────────────────────────────────────────────────────────────

\paragraph{Velocity tracking.}
\[
  r_\text{vel} = \exp\!\left(-\frac{\|v_{xy}^\text{cmd} - v_{xy}\|^2}{\sigma^2}\right),
  \quad \sigma^2 = 0.25
\]

\paragraph{Gait reward.}
\[
  r_\text{gait} = \sum_{i=1}^{2} \mathbf{1}\!\left[\text{stance}_i(t) = \text{contact}_i(t)\right]
\]
where $\text{stance}_i(t) = \mathbf{1}[\phi_i(t) < 0.55]$,
$\phi_i(t) = \left(\frac{t}{T_g} + \Delta\phi_i\right) \!\mod 1$,
$T_g = 0.8$\,s, $\Delta\phi_1 = 0$, $\Delta\phi_2 = 0.5$.

\paragraph{LiDAR proximity penalty.}
\[
  r_\text{lidar} = -\sum_{k=1}^{K} \frac{\max(0,\, d_\text{safe} - d_k)}{d_\text{safe}},
  \quad d_\text{safe} = 0.8\,\text{m},\; K = 248
\]

% ─────────────────────────────────────────────────────────────
\section{Observation Space Details}
\label{app:obs}
% ─────────────────────────────────────────────────────────────

\begin{center}
\begin{tabular}{lccc}
  \toprule
  \textbf{Obs.\ group} & \textbf{Modality} & \textbf{Per-step dim}
    & \textbf{With history ($T{=}5$)} \\
  \midrule
  Policy - no LiDAR & Proprioception & 96 & 480 \\
  Policy - with LiDAR & Proprio + LiDAR & 344 & 1{,}720 \\
  Critic - no LiDAR & Proprioception+ & 99 & 495 \\
  Critic - with LiDAR & Proprio+ + LiDAR & 347 & 1{,}735 \\
  \bottomrule
\end{tabular}
\end{center}

The critic's privileged observation adds ground-truth base linear velocity (3D)
to the policy observation to reduce value estimation variance during training
while keeping the actor blind to this privileged signal at deployment.

\end{document}
